\uuid{8pes}
\titre{Équation homogène/solution particulière}
\auteur{hanani ; mzari}
\organisation{exo7}
\licence{CC BY-NC-SA 4.0}
\datecreate{2020-01-20}
\niveau{L1}
\module{Analyse}
\chapitre{Equation différentielle}
\sousChapitre{Résolution d'équation différentielle du premier ordre}
\difficulte{3}
\theme{Equations du premier ordre, équation homogène, solution particulière}
\exo7qcmid{lille-2.equations-differentielles.007}

\contenu{
\begin{question}
\qclassification{225.01}
\qtags{motcle=équation homogène/solution particulière}
On considère l'équation différentielle $(E)$ : $\displaystyle \left(1+\cos ^2x\right)y'+y\sin (2x)=0$. Parmi les affirmations suivantes cocher celles qui sont vraies :
\begin{answers}
\bad{$\displaystyle A(x)=\ln (1+\cos ^2x)$ est une primitive sur $\Rr$ de $\displaystyle a(x)=\frac{\sin (2x)}{1+\cos ^2x}$.}
\good{Toute solution de $(E)$ vérifie $y'(0)=0$.}
\bad{$(E)$ n'admet pas de solution vérifiant $y(0)=0$.}
\good{La solution générale de $(E)$ est : $\displaystyle y=k+k\cos ^2x$, $k\in \Rr$.}
\end{answers}
\begin{explanations}
Toute solution $y$ de $(E)$ vérifie $y'(0)=0$ car ,en posant $x=0$ dans $(E)$, on obtient $y'(0)=0$. La relation $\sin(2x)=2\sin x \cos x$ montre que $a(x)$ est de la forme $\displaystyle \frac{-u'}{u}$ avec $u=1+\cos ^2x$. Donc toute primitive de $a$ est de la forme
$$A(x)=-\ln (1+\cos ^2x)+C,\quad C\in \Rr.$$
La solution générale de $(E)$ est $y=k(1+\cos ^2x)$, $k\in \Rr$.
\end{explanations}
\end{question}
}
